\documentclass[10pt,a4paper]{amsart}
\usepackage[margin=19mm]{geometry}
\usepackage{amsmath,amssymb,mathtools}
\usepackage[scr=rsfs,cal=euler]{mathalfa}
\usepackage{xfrac}
\usepackage[unicode,hidelinks,bookmarks=false]{hyperref}
\newcommand{\ie}{i.e.\ }
\newcommand{\cf}{cf.\ }
\newcommand{\resp}{respectively\ }
\newcommand{\Subsection}{\subsection}
\providecommand{\nobreakdash}{\nobreak\hbox{-}}
\setlength{\emergencystretch}{2em}
\multlinegap0pt
\usepackage{amssymb}
\usepackage[shortlabels]{enumitem}
\usepackage{mathtools}
\newtheorem{theorem}{Theorem}
\newcommand{\End}{{\rm End}}
\newcommand{\Q}{{\mathbb Q}}
\newcommand{\Z}{{\mathbb Z}}
\newcommand{\eps}{{\varepsilon}}
\newcommand{\tens}{\otimes}
\title{Hybrid conjecture in a mixed Shimura variety}
\author{Rodolphe Richard, Andrei Yafaev}
\date{Source: arXiv:2604.23376v1 (CC BY 4.0)}
\begin{document}
\maketitle

\noindent\emph{Source and licence.} Hybrid conjecture in a mixed Shimura variety. Version arXiv:2604.23376v1 (CC BY 4.0), distributed by arXiv under the Creative Commons Attribution 4.0 International licence, http://creativecommons.org/licenses/by/4.0/, as recorded in the arXiv licence metadata for this exact version and shown on the abstract page https://arxiv.org/abs/2604.23376v1. Full licence notice: \texttt{licenses/arxiv-2604.23376-CC-BY-4.0.txt}.\par\smallskip

\noindent\textit{Editorial scope.} Counted unit: the article's theorem linear (source0.tex lines 3049--3065). The article's complete original proof of it is delivered with it (source0.tex lines 3067--3112, one \texttt{proof} environment of 46 lines). No mathematics is written, repaired or re-derived: the statement and the proof are the article's own lines, in the article's own notation, and no retained line is substituted or re-flowed.  No further source-local context is needed: every cross-reference of every retained line resolves inside the two delivered spans.  Nothing was omitted from any retained span, so the package carries no omission record.  No retained line contains a citation, so the wrapper carries no bibliography and no bibliography entry.  No retained line contains a figure environment, an image-inclusion command or drawing code, so no image is delivered and the package declares no asset dependency.  Editorial formatting only, in addition: an AMS \texttt{amsart} preamble that restates the article's own declarations and macros used by the retained lines, namely the environments \texttt{theorem} on separate counters, together with the standard packages those lines need.  The retained environments print in this document as Theorem 1 in order of appearance, whereas the article prints the same items with the numbering of its own full document.  The complete native source of the article as distributed in the arXiv e-print for this exact version is bundled unchanged as \texttt{sources/199159/source0.tex}.
\begin{theorem}\label{thm:split linear}
Let~$L/\Q_\ell$ be a finite extension, let~$k,n_1,\ldots,n_k\in\Z_{\geq1}$ and let~$V_1,\ldots,V_k$ be finite dimensional $L$-vector spaces.
Let~$W\leq V= \bigoplus_{i=1}^k L^{n_i}\tens V_i$ be an~$O_L$-submodule.  We consider the action
\[
E:=\prod_{i=1}^k \End(O_L^{n_i})\tens Id_{V_i}\leq \prod_{i=1}^k\End(L^{n_i})\tens \End(V_i)\leq \End(V).
\]
Let~$\lambda\in L$ be such that
\[
\lambda \cdot E\cdot W\leq W.
\]
Then there are~$O_L$-modules~$\Lambda_i\leq V_i$ such that
\[
\lambda\cdot\bigoplus_{i=1}^k O_L^{n_i}\tens_{O_L} \Lambda_i\leq W
\text{ and }
\lambda\cdot W\leq \bigoplus_{i=1}^k O_L^{n_i}\tens_{O_L}\Lambda_i.
\]
\end{theorem}

\begin{proof}[Proof of theorem~\ref{thm:split linear}]
Let~$e_{i,1},\ldots,e_{i,n_i}$ be the standard basis of~${L}^{n_i}$, let
\[
\eps_{i,(m,m')}:= e_{i,m}\times {}^te_{i,m'}\in \End(L^{n_i}):(\lambda_n)_{1\leq n\leq n_i}\mapsto (\delta_{m,n} \lambda_{m'})_{1\leq n\leq n_i}
\]
be the standard elementary matrices and let~
\[
\pi_{i,m}=\eps_{i,(m,m)}\in \End(L^{n_i}):(\lambda_n)_{1\leq n\leq n_i}\mapsto (\delta_{m,n} \lambda_n)_{1\leq n\leq n_i}
\]
be the~$m$-th standard projector.

We can uniquely decompose any~$v\in V$ as
\[
v=\bigoplus_{i=1}^k \sum^{n_i}_{m=1} e_{i,m}\tens v_{i,m}.
\]
Let~$\Lambda_{i,m}:=\{v\in V_i|e_{i,m}\tens v\in W\}$. If~$v\in \Lambda_{i,m}$,
then
\[
e_{i,m'}\tens \lambda \cdot v = \lambda \eps_{i,(m',m)}(e_{i,m}\tens v)\in \lambda\cdot E(W)\leq W.
\]
Therefore~$\lambda\cdot \Lambda_{i,m}\leq \Lambda_{i,m'}$. We deduce
\[
\Lambda_i:=\sum_{i=1}^{m}\Lambda_{i,m'}\geq \bigcap_{i=1}^{n_i} \Lambda_{i,m}\geq \lambda\cdot \Lambda_i.
\]
By construction
\begin{equation}\label{eq:356}
\bigoplus_{i=1}^k O_L^{n_i}\tens \lambda\cdot \Lambda_i
\leq 
\bigoplus_{i=1}^k \sum^{n_i}_{m=1} e_{i,m}\tens \Lambda_{i,m}
\leq W.
\end{equation}
For~$w\in W$, we have
\(
e_i\tens \lambda w_{i,m}=\lambda \pi_{i,m}(w)\in W.
\)
Therefore~$\lambda\cdot w_{i,m}\in \Lambda_{i,m}$ and
\[
\lambda\cdot w=\sum_{1\leq i\leq k,1\leq m\leq n_i}e_{i,m}\tens \lambda\cdot w_{i,m}\in 
\bigoplus_{i=1}^k \sum^{n_i}_{m=1} e_{i,m}\tens \Lambda_{i,m}\leq \bigoplus_{i=1}^k O_L^{n_i}\tens \Lambda_i.
\]
As~$w\in W$ is arbitrary, we conclude
\begin{equation}\label{eq:123}
\lambda\cdot W\leq \bigoplus_{i=1}^k O_L^{n_i}\tens \Lambda_i.
\end{equation}
The Theorem follows from~\eqref{eq:356} and~\eqref{eq:123}.
\end{proof}

\end{document}
